\documentclass[10pt,a4paper]{amsart}
\usepackage[margin=19mm]{geometry}
\usepackage{amsmath,amssymb,mathtools}
\usepackage[scr=rsfs,cal=euler]{mathalfa}
\usepackage{xfrac}
\usepackage[unicode,hidelinks,bookmarks=false]{hyperref}
\newcommand{\ie}{i.e.\ }
\newcommand{\cf}{cf.\ }
\newcommand{\resp}{respectively\ }
\newcommand{\Subsection}{\subsection}
\providecommand{\nobreakdash}{\nobreak\hbox{-}}
\setlength{\emergencystretch}{2em}
\multlinegap0pt
\usepackage{amssymb}
\newtheorem{proposition}{Proposition}
\newcommand{\Coind}{\mathrm{CoInd}}
\newcommand{\Nm}{\mathrm{Nm}}
\newcommand{\Spec}{\mathrm{Spec}}
\newcommand{\Tr}{\mathrm{Tr}}
\newcommand{\ghost}{\Gamma}
\newcommand{\p}{\mathfrak{p}}
\newcommand{\q}{\mathfrak{q}}
\renewcommand{\subset}{\subseteq}
\title{The subgroup stratification of Nakaoka spectra}
\author{Noah Wisdom}
\date{Source: arXiv:2508.09360v2 (CC BY 4.0)}
\begin{document}
\maketitle

\noindent\emph{Source and licence.} The subgroup stratification of Nakaoka spectra. Version arXiv:2508.09360v2 (CC BY 4.0), distributed by arXiv under the Creative Commons Attribution 4.0 International licence, http://creativecommons.org/licenses/by/4.0/, as recorded in the arXiv licence metadata for this exact version and shown on the abstract page https://arxiv.org/abs/2508.09360v2. Full licence notice: \texttt{licenses/arxiv-2508.09360-CC-BY-4.0.txt}.\par\smallskip

\noindent\textit{Editorial scope.} Counted unit: the article's proposition prop:stratum-of-ghost (source0.tex lines 683--690). The article's complete original proof of it is delivered with it (source0.tex lines 692--706, one \texttt{proof} environment of 15 lines). No mathematics is written, repaired or re-derived: the statement and the proof are the article's own lines, in the article's own notation, and no retained line is substituted or re-flowed.  No further source-local context is needed: every cross-reference of every retained line resolves inside the two delivered spans.  Nothing was omitted from any retained span, so the package carries no omission record.  No retained line contains a citation, so the wrapper carries no bibliography and no bibliography entry.  No retained line contains a figure environment, an image-inclusion command or drawing code, so no image is delivered and the package declares no asset dependency.  Editorial formatting only, in addition: an AMS \texttt{amsart} preamble that restates the article's own declarations and macros used by the retained lines, namely the environments \texttt{proposition} on separate counters, together with the standard packages those lines need.  The retained environments print in this document as Proposition 1 in order of appearance, whereas the article prints the same items with the numbering of its own full document.  The complete native source of the article as distributed in the arXiv e-print for this exact version is bundled unchanged as \texttt{sources/183157/source0.tex}.
\begin{proposition}\label{prop:stratum-of-ghost}
    Let $R$ be any nonzero $C_p$-Tambara functor and $\ghost(R)$ its ghost. 
    \begin{enumerate}
        \item A prime ideal $(\p;\q)$ of $\ghost(R)$ belongs to the $e$th stratum of $\Spec(\ghost(R))$ if and only if it is equal to $(\p;R(C_p/C_p)/\Tr)$ for $\p$ a prime ideal of $R(C_p/e)$.
        \item If $\Tr$ is not surjective then the $e$th stratum of $\Spec(\ghost(R))$ is not open.
        \item If $R(C_p/e)^{C_p}$ has finitely many minimal prime ideals (e.g. is Noetherian or an integral domain), the $e$th stratum of $\Spec(\ghost(R))$ is closed.
    \end{enumerate}
\end{proposition}

\begin{proof}
    The first statement follows from the observation that the map $\ghost(R) \rightarrow \Coind_e^{C_p} R(C_p/e)$ is given in level $C_p/C_p$ by the restriction 
    \[ 
        \ghost(R)(C_p/C_p) \rightarrow \ghost(R)(C_p/e) = R(C_p/e) . 
    \] 
    This restriction is the projection 
    \[ 
        \ghost(R)(C_p/C_p) = R(C_p/e)^{C_p} \times R(C_p/C_p)/\Tr \rightarrow R(C_p/e)^{C_p} . 
    \] 
    The preimage of an arbitrary prime ideal $\Coind_e^{C_p} \p$ of $\Coind_e^{C_p} R(C_p/e)$ is therefore $\bigcap_{g \in C_p} g \p$ in the $C_p/e$-level and $\p^{C_p} \times R(C_p/e)/\Tr$ in the $C_p/C_p$-level.

    For the second statement, let $\p$ be any prime ideal of $R(C_p/C_p)/\Tr$. Then $(\Nm^{-1}\p;\p)$ is not in the $e$th stratum, but the closure of $(\Nm^{-1}\p;\p)$ contains the point $(\Nm^{-1}\p;R(C_p/C_p)/\Tr)$ which is in the $e$th stratum. Thus the complement of the $e$th stratum is not closed.

    Finally, let $\p_i$ be the finite list of minimal prime ideals of $R(C_p/e)^{C_p}$. By the first paragraph, every prime ideal in the $e$th stratum has the form $(\q;R(C_p/C_p)/\Tr)$. Now for some $i$ we have $\p_i \subset \q$, whence $(\p_i;R(C_p/C_p)/\Tr) \subset (\q;R(C_p/C_p)/\Tr)$. Therefore the $e$th stratum is the finite union of the basic closed subsets $V((\p_i);R(C_p/C_p)/\Tr)$.
\end{proof}

\end{document}
